\section{Problem Definition and Analytics Objective}\label{sec:problem}

\textbf{Context.} India's analytics job market changes faster than curricula. Universities, skilling
providers and CSR programmes must decide what to teach with limited budget, trainers and seats,
but lack evidence on three questions: which skills the market \emph{demands and pays for}, which
skills are associated with \emph{early-career growth}, and which personal traits are associated
with \emph{success among senior} data scientists. The SAS-provided data covers all three: job
postings (Analytics Jobs), company-level openings and salaries (DataScience Jobs), and two
outcome-labelled employee datasets (JDS Skill Traits for juniors, SDS Personality Traits for
seniors).

\begin{center}
\fcolorbox{sasblue}{sasblue!6}{\parbox{0.94\linewidth}{\textbf{Analytics objective.}
To identify the technical skills and personality traits most strongly associated with
labour-market demand, pay and on-the-job success for data-science professionals, and to
translate this evidence into a budget-constrained, explainable training plan that tells an
institution \emph{what to teach first}.}}
\end{center}

\begin{table}[H]\centering\small
\caption{Research questions, data and methods.}\label{tab:rq}
\begin{tabularx}{\linewidth}{@{}L{2.0cm}Y L{3.4cm} L{3.9cm}@{}}
\toprule
\textbf{RQ} & \textbf{Question} & \textbf{Data} & \textbf{Methods}\\
\midrule
RQ1 Demand \& pay & Which skills and roles are most demanded, and which are associated with
higher salary bands, beyond experience? & Analytics Jobs (14,840 cleaned postings); DataScience Jobs
(1,602 rows, 93,005 openings) & Cleaning, skill normalisation (ESCO/O*NET), Wilson CIs, ordinal
logistic regression, FDR control, association rules \\
\addlinespace
RQ2 Junior growth & Which of five measured skill areas are associated with a high vs.\ low salary
hike? & JDS Skill Traits (139) & Mann--Whitney U, effect sizes, regularised logistic regression,
shallow tree, repeated CV, permutation test \\
\addlinespace
RQ3 Senior success & Which Big Five traits are associated with high vs.\ low success? & SDS
Personality Traits (161) & As RQ2 \\
\addlinespace
RQ4 Decision & Given RQ1--RQ3, which training interventions should an institution prioritise under
budget and capacity limits? & Evidence from RQ1--RQ3; documented cost assumptions & Integer
optimisation (CP-SAT), sensitivity, Pareto frontier, robustness \\
\bottomrule
\end{tabularx}
\end{table}

\textbf{Stakeholders.} The primary decision-maker is a training institution (university,
skilling provider or CSR programme). Secondary stakeholders are students choosing what to learn,
employers and HR teams who write job descriptions and assess staff, and policy makers allocating
skilling funds (Section~\ref{sec:implications}).

\textbf{Scope.} Only the SAS-provided datasets are analysed; ESCO and O*NET
\citep{esco2024,onet2025} serve solely as a skill vocabulary. Findings are associations, not
causal effects. Analytics Jobs contains non-data roles too; we keep them for comparison and analyse
the 5,138 data-role postings separately (Section~\ref{sec:families}).

\textbf{Success criteria.} Each RQ is answered with a quantified finding; predictive models are
judged against a majority-class baseline with cross-validated accuracy and AUC; the RQ4 plan states
which constraints bind and what extra resources would change.
